%%%%%%%%%%%%%%%%%%%%%%%%%%%%%%%%%%%%%%%%%%%%%%%%%%%%%%%%%%%%%%%%%%%%%%
% Lớp 8 | Chương II. Hằng đẳng thức đáng nhớ và ứng dụng | Tổng và hiệu hai lập phương
% Dạng 8D2?8-2: Rút gọn, tính giá trị biểu thức
%%%%%%%%%%%%%%%%%%%%%%%%%%%%%%%%%%%%%%%%%%%%%%%%%%%%%%%%%%%%%%%%%%%%%%


% nguồn: De_AI_Toan8_cuoi_hoc_ki_i_1 (nhập 21/11/2026)
\begin{bt}%[8D2H8-2]%[ma:fb83fee4]
	Tính nhanh các giá trị sau bằng cách áp dụng hằng đẳng thức:
a) $75^2 - 25^2$;
b) $86^2 - 14^2$.
	\loigiai{
		a) $75^2 - 25^2 = (75 - 25)(75 + 25) = 50 \cdot 100 = 5000$. \\
b) $86^2 - 14^2 = (86 - 14)(86 + 14) = 72 \cdot 100 = 7200$.
	}
\end{bt}

% nguồn: De_AI_Toan8_cuoi_hoc_ki_i_1 (nhập 21/11/2026)
\begin{bt}%[8D2H8-2]%[ma:b0b13e54]
	Tính nhanh bằng hằng đẳng thức:
a) $101^2$;
b) $99^2$.
	\loigiai{
		a) $101^2 = (100 + 1)^2 = 100^2 + 2 \cdot 100 \cdot 1 + 1^2 = 10000 + 200 + 1 = 10201$. \\
b) $99^2 = (100 - 1)^2 = 100^2 - 2 \cdot 100 \cdot 1 + 1^2 = 10000 - 200 + 1 = 9801$.
	}
\end{bt}

% nguồn: De_AI_Toan8_cuoi_hoc_ki_i_1 (nhập 21/11/2026)
\begin{bt}%[8D2H8-2]%[ma:42b6941a]
	Tính nhanh các tích sau bằng hằng đẳng thức hiệu hai bình phương:
a) $48 \cdot 52$;
b) $97 \cdot 103$.
	\loigiai{
		a) $48 \cdot 52 = (50 - 2)(50 + 2) = 50^2 - 2^2 = 2500 - 4 = 2496$. \\
b) $97 \cdot 103 = (100 - 3)(100 + 3) = 100^2 - 3^2 = 10000 - 9 = 9991$.
	}
\end{bt}

% nguồn: De_AI_Toan8_cuoi_hoc_ki_i_1 (nhập 21/11/2026)
\begin{bt}%[8D2H8-2]%[ma:886b4e5d]
	Tính nhanh giá trị các biểu thức số học sau:
a) $37^2 + 2 \cdot 37 \cdot 13 + 13^2$;
b) $65^2 - 2 \cdot 65 \cdot 15 + 15^2$.
	\loigiai{
		a) $37^2 + 2 \cdot 37 \cdot 13 + 13^2 = (37 + 13)^2 = 50^2 = 2500$. \\
b) $65^2 - 2 \cdot 65 \cdot 15 + 15^2 = (65 - 15)^2 = 50^2 = 2500$.
	}
\end{bt}

% nguồn: De_AI_Toan8_cuoi_hoc_ki_i_1 (nhập 21/11/2026)
\begin{bt}%[8D2H8-2]%[ma:fc524888]
	Tính nhanh bằng hằng đẳng thức lập phương của một tổng hoặc một hiệu:
a) $19^3$;
b) $51^3$.
	\loigiai{
		a) $19^3 = (20 - 1)^3 = 20^3 - 3 \cdot 20^2 \cdot 1 + 3 \cdot 20 \cdot 1^2 - 1^3 = 8000 - 1200 + 60 - 1 = 6859$. \\
b) $51^3 = (50 + 1)^3 = 50^3 + 3 \cdot 50^2 \cdot 1 + 3 \cdot 50 \cdot 1^2 + 1^3 = 125000 + 7500 + 150 + 1 = 132651$.
	}
\end{bt}

% nguồn: De_AI_Toan8_cuoi_hoc_ki_i_1 (nhập 21/11/2026)
\begin{bt}%[8D2H8-2]%[ma:8a57f2f5]
	Tính nhanh bằng hằng đẳng thức:
a) $102^3$;
b) $98^3$.
	\loigiai{
		a) $102^3 = (100 + 2)^3 = 100^3 + 3 \cdot 100^2 \cdot 2 + 3 \cdot 100 \cdot 2^2 + 2^3 = 1000000 + 60000 + 1200 + 8 = 1061208$. \\
b) $98^3 = (100 - 2)^3 = 100^3 - 3 \cdot 100^2 \cdot 2 + 3 \cdot 100 \cdot 2^2 - 2^3 = 1000000 - 60000 + 1200 - 8 = 941192$.
	}
\end{bt}

% nguồn: De_AI_Toan8_cuoi_hoc_ki_i_1 (nhập 21/11/2026)
\begin{bt}%[8D2H8-2]%[ma:73ae28e0]
	Tính nhanh giá trị của biểu thức $A = \dfrac{73^3 + 27^3}{100} - 73 \cdot 27$.
	\loigiai{
		Áp dụng hằng đẳng thức tổng hai lập phương $a^3 + b^3 = (a + b)(a^2 - ab + b^2)$: \\
$73^3 + 27^3 = (73 + 27)(73^2 - 73 \cdot 27 + 27^2) = 100(73^2 - 73 \cdot 27 + 27^2)$. \\
Do đó $A = (73^2 - 73 \cdot 27 + 27^2) - 73 \cdot 27 = 73^2 - 2 \cdot 73 \cdot 27 + 27^2 = (73 - 27)^2 = 46^2 = 2116$.
	}
\end{bt}

% nguồn: De_AI_Toan8_cuoi_hoc_ki_i_1 (nhập 21/11/2026)
\begin{bt}%[8D2H8-2]%[ma:a745150d]
	Tính nhanh giá trị của biểu thức $B = \dfrac{64^3 - 14^3}{50} + 64 \cdot 14$.
	\loigiai{
		Áp dụng hằng đẳng thức hiệu hai lập phương $a^3 - b^3 = (a - b)(a^2 + ab + b^2)$: \\
$64^3 - 14^3 = (64 - 14)(64^2 + 64 \cdot 14 + 14^2) = 50(64^2 + 64 \cdot 14 + 14^2)$. \\
Do đó $B = (64^2 + 64 \cdot 14 + 14^2) + 64 \cdot 14 = 64^2 + 2 \cdot 64 \cdot 14 + 14^2 = (64 + 14)^2 = 78^2 = 6084$.
	}
\end{bt}

% nguồn: De_AI_Toan8_cuoi_hoc_ki_i_1 (nhập 21/11/2026)
\begin{bt}%[8D2H8-2]%[ma:ea1ab919]
	Tính nhanh giá trị của phân số $C = \dfrac{85^2 - 15^2}{55^2 - 45^2}$.
	\loigiai{
		Áp dụng hằng đẳng thức hiệu hai bình phương cho cả tử và mẫu: \\
Tử số: $85^2 - 15^2 = (85 - 15)(85 + 15) = 70 \cdot 100 = 7000$. \\
Mẫu số: $55^2 - 45^2 = (55 - 45)(55 + 45) = 10 \cdot 100 = 1000$. \\
Suy ra $C = \dfrac{7000}{1000} = 7$.
	}
\end{bt}

% nguồn: De_AI_Toan8_cuoi_hoc_ki_i_1 (nhập 21/11/2026)
\begin{bt}%[8D2H8-2]%[ma:3432849e]
	Cho $x + y = 7$ và $xy = 10$. Tính giá trị của các biểu thức đối xứng sau:
a) $x^2 + y^2$;
b) $(x - y)^2$.
	\loigiai{
		a) Ta có $x^2 + y^2 = (x + y)^2 - 2xy = 7^2 - 2 \cdot 10 = 49 - 20 = 29$. \\
b) Ta có $(x - y)^2 = (x + y)^2 - 4xy = 7^2 - 4 \cdot 10 = 49 - 40 = 9$.
	}
\end{bt}

% nguồn: De_AI_Toan8_cuoi_hoc_ki_i_1 (nhập 21/11/2026)
\begin{bt}%[8D2H8-2]%[ma:46b6a745]
	Cho $x - y = 4$ và $xy = 5$. Tính giá trị của các biểu thức:
a) $x^2 + y^2$;
b) $(x + y)^2$.
	\loigiai{
		a) Ta có $x^2 + y^2 = (x - y)^2 + 2xy = 4^2 + 2 \cdot 5 = 16 + 10 = 26$. \\
b) Ta có $(x + y)^2 = (x - y)^2 + 4xy = 4^2 + 4 \cdot 5 = 16 + 20 = 36$.
	}
\end{bt}

% nguồn: De_AI_Toan8_cuoi_hoc_ki_i_1 (nhập 21/11/2026)
\begin{bt}%[8D2H8-2]%[ma:f49493f6]
	Cho $x + y = 6$ và $xy = 8$. Tính giá trị của biểu thức $x^3 + y^3$.
	\loigiai{
		Ta có hằng đẳng thức $x^3 + y^3 = (x + y)^3 - 3xy(x + y)$. \\
Thay $x + y = 6$ và $xy = 8$ vào biểu thức, ta được: \\
$x^3 + y^3 = 6^3 - 3 \cdot 8 \cdot 6 = 216 - 144 = 72$.
	}
\end{bt}

% nguồn: De_AI_Toan8_cuoi_hoc_ki_i_1 (nhập 21/11/2026)
\begin{bt}%[8D2H8-2]%[ma:73c032e6]
	Cho $x - y = 3$ và $xy = 10$. Tính giá trị của biểu thức $x^3 - y^3$.
	\loigiai{
		Ta có hằng đẳng thức $x^3 - y^3 = (x - y)^3 + 3xy(x - y)$. \\
Thay $x - y = 3$ và $xy = 10$ vào biểu thức, ta được: \\
$x^3 - y^3 = 3^3 + 3 \cdot 10 \cdot 3 = 27 + 90 = 117$.
	}
\end{bt}

% nguồn: De_AI_Toan8_cuoi_hoc_ki_i_1 (nhập 21/11/2026)
\begin{bt}%[8D2H8-2]%[ma:d8b529b8]
	Cho $x + y = 5$ và $x^2 + y^2 = 17$. Tính:
a) $xy$;
b) $x^3 + y^3$.
	\loigiai{
		a) Từ $(x + y)^2 = x^2 + 2xy + y^2$, suy ra $2xy = (x + y)^2 - (x^2 + y^2) = 5^2 - 17 = 8$, nên $xy = 4$. \\
b) Áp dụng hằng đẳng thức tổng hai lập phương: \\
$x^3 + y^3 = (x + y)(x^2 - xy + y^2) = 5 \cdot (17 - 4) = 5 \cdot 13 = 65$.
	}
\end{bt}

% nguồn: De_AI_Toan8_cuoi_hoc_ki_i_1 (nhập 21/11/2026)
\begin{bt}%[8D2H8-2]%[ma:25e195aa]
	Cho $x - y = 2$ và $x^2 + y^2 = 20$. Tính:
a) $xy$;
b) $x^3 - y^3$.
	\loigiai{
		a) Từ $(x - y)^2 = x^2 - 2xy + y^2$, suy ra $2xy = (x^2 + y^2) - (x - y)^2 = 20 - 2^2 = 16$, nên $xy = 8$. \\
b) Áp dụng hằng đẳng thức hiệu hai lập phương: \\
$x^3 - y^3 = (x - y)(x^2 + xy + y^2) = 2 \cdot (20 + 8) = 2 \cdot 28 = 56$.
	}
\end{bt}

% nguồn: De_AI_Toan8_cuoi_hoc_ki_i_1 (nhập 21/11/2026)
\begin{bt}%[8D2H8-2]%[ma:343df226]
	Cho số thực $x \ne 0$ thỏa mãn $x + \dfrac{1}{x} = 4$. Tính giá trị của các biểu thức:
a) $x^2 + \dfrac{1}{x^2}$;
b) $x^3 + \dfrac{1}{x^3}$.
	\loigiai{
		a) Ta có $\left(x + \dfrac{1}{x}\right)^2 = x^2 + 2 \cdot x \cdot \dfrac{1}{x} + \dfrac{1}{x^2} = x^2 + \dfrac{1}{x^2} + 2$. \\
Suy ra $x^2 + \dfrac{1}{x^2} = 4^2 - 2 = 16 - 2 = 14$. \\
b) Ta có $x^3 + \dfrac{1}{x^3} = \left(x + \dfrac{1}{x}\right)^3 - 3 \cdot x \cdot \dfrac{1}{x} \left(x + \dfrac{1}{x}\right) = 4^3 - 3 \cdot 4 = 64 - 12 = 52$.
	}
\end{bt}

% nguồn: De_AI_Toan8_cuoi_hoc_ki_i_1 (nhập 21/11/2026)
% nguồn: chuyen_dang (nhập 22/11/2026)
\begin{bt}%[8D2H8-2]%[ma:e8b7ace8]
	Cho $x + y = 3$ và $xy = -2$. Tính giá trị của biểu thức $M = (x + 2)(y + 2)$.
	\loigiai{
		Ta có $M = (x + 2)(y + 2) = xy + 2x + 2y + 4 = xy + 2(x + y) + 4$. \\
Thay $x + y = 3$ và $xy = -2$ vào $M$, ta được $M = -2 + 2 \cdot 3 + 4 = 8$.
	}
\end{bt}

% nguồn: chuyen_dang (nhập 22/11/2026)
\begin{bt}%[8D2H8-2]%[ma:77122d0a]
	Cho $x - y = 5$ và $xy = 6$. Tính giá trị của biểu thức $N = (x - 3)(y + 3)$.
	\loigiai{
		Ta có $N = (x - 3)(y + 3) = xy + 3x - 3y - 9 = xy + 3(x - y) - 9$. \\
Thay $x - y = 5$ và $xy = 6$ vào $N$, ta được $N = 6 + 3 \cdot 5 - 9 = 12$.
	}
\end{bt}

% nguồn: chuyen_dang (nhập 22/11/2026)
\begin{bt}%[8D2H8-2]%[ma:430fc03b]
	Tính nhanh giá trị của biểu thức $D = 103^3 - 9 \cdot 103^2 + 27 \cdot 103 - 27$.
	\loigiai{
		Ta biến đổi biểu thức $D$ về dạng lập phương của một hiệu:
		\begin{align*} 
			D &= 103^3 - 3 \cdot 103^2 \cdot 3 + 3 \cdot 103 \cdot 3^2 - 3^3 \\
			&= (103-3)^3 \\
			&= 100^3 \\
			&= 1\,000\,000.
		\end{align*}
	}
\end{bt}

% nguồn: chuyen_dang (nhập 22/11/2026)
\begin{bt}%[8D2N8-2]%[ma:3a65d0cb]
	Tính nhanh bằng cách áp dụng hằng đẳng thức:
a) $98^2 + 4 \cdot 98 + 4$;
b) $105^2 - 10 \cdot 105 + 25$.
	\loigiai{
		a) $98^2 + 4 \cdot 98 + 4 = 98^2 + 2 \cdot 98 \cdot 2 + 2^2 = (98 + 2)^2 = 100^2 = 10000$. \\
b) $105^2 - 10 \cdot 105 + 25 = 105^2 - 2 \cdot 105 \cdot 5 + 5^2 = (105 - 5)^2 = 100^2 = 10000$.
	}
\end{bt}

% nguồn: chuyen_dang (nhập 22/11/2026)
\begin{bt}%[8D2H8-2]%[ma:8744438a]
	Tính nhanh giá trị của biểu thức $C = 28^3 + 3 \cdot 28^2 \cdot 2 + 3 \cdot 28 \cdot 2^2 + 2^3$.
	\loigiai{
		Ta thấy biểu thức $C$ có dạng khai triển của lập phương một tổng:
		\begin{align*} 
			C &= 28^3 + 3 \cdot 28^2 \cdot 2 + 3 \cdot 28 \cdot 2^2 + 2^3 \\
			&= (28+2)^3 \\
			&= 30^3 \\
			&= 27\,000.
		\end{align*}
	}
\end{bt}
